% ===========================================================================
% M.Tech Project Report -- Swaral Gaur
% Supervisor: Mrs. Naina Yadav
%
% BEFORE COMPILING, FILL IN:
%   1. \myrollno  -- your roll number (search for "XXXXXXXX" below)
%   2. \supervisordesignation -- verify Mrs. Naina Yadav's exact designation
%      (currently set to "Assistant Professor")
%
% Compiles with pdfLaTeX on Overleaf. No external images required.
% ===========================================================================

\documentclass[12pt,a4paper,oneside]{report}

\usepackage[T1]{fontenc}
\usepackage[utf8]{inputenc}
\usepackage[a4paper,left=1.5in,right=1in,top=1in,bottom=1in]{geometry}
\usepackage{mathptmx}   % Times for both text and mathematics
\usepackage{setspace}
\usepackage{amsmath,amssymb}
\usepackage{graphicx}
\usepackage{array}
\usepackage{booktabs}
\usepackage{tabularx}
\usepackage{longtable}
\usepackage{enumitem}
\usepackage{titlesec}
\usepackage{tocloft}
\usepackage[hidelinks]{hyperref}

\onehalfspacing
\setlength{\parindent}{0pt}
\setlength{\parskip}{6pt}

% ---- Chapter heading format:  "CHAPTER - 1"  then title, both centred ----
\titleformat{\chapter}[display]
  {\normalfont\large\bfseries\centering}
  {CHAPTER - \thechapter}
  {8pt}
  {\large\MakeUppercase}
\titlespacing*{\chapter}{0pt}{-20pt}{24pt}

% Unnumbered chapters (e.g. REFERENCES) use the same centred style, no label
\titleformat{name=\chapter,numberless}[display]
  {\normalfont\large\bfseries\centering}
  {}
  {0pt}
  {\large\MakeUppercase}

\titleformat{\section}
  {\normalfont\normalsize\bfseries}{\thesection}{1em}{}

% ---- Personal details ----
\newcommand{\myname}{SWARAL GAUR}
\newcommand{\mynamemixed}{Swaral Gaur}
\newcommand{\myrollno}{XXXXXXXX}          % <-- FILL IN YOUR ROLL NUMBER
\newcommand{\supervisorname}{Mrs. Naina Yadav}
\newcommand{\supervisordesignation}{Assistant Professor}  % <-- VERIFY
\newcommand{\reporttitle}{Reliability-Weighted Prompt Ensembling for Cross-Domain
Session Recommendation Using Large Language Models}
\newcommand{\reporttitleshort}{Reliability-Weighted Prompt Ensembling for
Cross-Domain Session Recommendation Using Large Language Models}

\begin{document}

% ======================================================= TITLE PAGE
\thispagestyle{empty}
\begin{center}
\vspace*{0.3in}
{\large\bfseries RELIABILITY-WEIGHTED PROMPT ENSEMBLING\\[4pt]
FOR CROSS-DOMAIN SESSION RECOMMENDATION\\[4pt]
USING LARGE LANGUAGE MODELS\par}

\vspace{0.5in}
{\normalsize A Project Report\par}
\vspace{6pt}
{\small submitted in partial fulfilment of the requirements for the award of the degree of\par}

\vspace{0.25in}
{\large\bfseries MASTER OF TECHNOLOGY\par}
\vspace{6pt}
{\normalsize in\par}
\vspace{6pt}
{\large\bfseries COMPUTER SCIENCE AND ENGINEERING\par}

\vspace{0.45in}
{\normalsize Under the supervision of\par}
\vspace{6pt}
{\large\bfseries \supervisorname\par}
{\normalsize \supervisordesignation\par}

\vspace{0.35in}
{\normalsize Submitted by\par}
\vspace{6pt}
{\large\bfseries \myname\par}
{\normalsize Roll No. \myrollno\par}

\vfill
{\normalsize\bfseries DR B R AMBEDKAR NATIONAL INSTITUTE OF TECHNOLOGY\par}
{\normalsize DEPARTMENT OF COMPUTER SCIENCE AND ENGINEERING\par}
{\normalsize JALANDHAR, PUNJAB -- 144008, INDIA\par}
\vspace{8pt}
{\normalsize\bfseries 2026\par}
\end{center}

\clearpage
\pagenumbering{roman}
\setcounter{page}{1}

% ======================================================= CERTIFICATE
\thispagestyle{plain}
\addcontentsline{toc}{chapter}{Certificate}
\begin{center}
{\large\bfseries CERTIFICATE}
\end{center}
\vspace{0.2in}

This is to certify that the project report entitled ``\reporttitle'' submitted by
\mynamemixed{} (Roll No. \myrollno) to the Department of Computer Science and
Engineering, Dr B R Ambedkar National Institute of Technology, Jalandhar, in
partial fulfilment of the requirements for the award of the degree of Master of
Technology in Computer Science and Engineering, is a record of bona fide work
carried out by him under my supervision and guidance.

To the best of my knowledge, the matter presented in this report has not been
submitted elsewhere for the award of any other degree or diploma.

\vspace{0.8in}

\noindent
\begin{tabularx}{\textwidth}{@{}Xr@{}}
Date: & \supervisorname\\[4pt]
Place: Jalandhar & \supervisordesignation\\[4pt]
 & Department of Computer Science and Engineering\\
\end{tabularx}

\clearpage

% ======================================================= DECLARATION
\thispagestyle{plain}
\addcontentsline{toc}{chapter}{Declaration}
\begin{center}
{\large\bfseries DECLARATION}
\end{center}
\vspace{0.2in}

I hereby declare that the work presented in this project report entitled
``\reporttitle'' is my own work carried out under the supervision of
\supervisorname, \supervisordesignation, Department of Computer Science and
Engineering, Dr B R Ambedkar National Institute of Technology, Jalandhar.

I further declare that the material obtained from other sources has been duly
acknowledged and cited in the report. This work has not been submitted, in part
or in full, to any other institute or university for the award of any degree or
diploma.

\vspace{1.0in}

\noindent
\begin{flushright}
\mynamemixed\\
Roll No. \myrollno\\
M.Tech (CSE)
\end{flushright}

\clearpage

% ======================================================= ACKNOWLEDGEMENT
\thispagestyle{plain}
\addcontentsline{toc}{chapter}{Acknowledgement}
\begin{center}
{\large\bfseries ACKNOWLEDGEMENT}
\end{center}
\vspace{0.2in}

I express my sincere gratitude to my supervisor, \supervisorname,
\supervisordesignation, Department of Computer Science and Engineering, for her
valuable guidance, constant encouragement and patient supervision throughout
this work. Her insistence on controlled experimental design, and on
distinguishing measurement from mitigation, shaped the direction of this study.

I am thankful to the Head of the Department and to all faculty members of the
Department of Computer Science and Engineering, Dr B R Ambedkar National
Institute of Technology, Jalandhar, for providing the facilities and academic
environment necessary to carry out this work.

I also acknowledge the GroupLens research group for the MovieLens collection and
the McAuley Laboratory at the University of California San Diego for the Amazon
Reviews 2023 collection, both of which are released publicly and without which
this study would not have been possible. I further acknowledge the developers of
the Qwen model family and of the Ollama runtime, whose openly available tools
made local, reproducible experimentation feasible on modest hardware.

Finally, I thank my family and friends for their continuous support and
motivation.

\vspace{0.6in}

\noindent
\begin{flushright}
\mynamemixed\\
Roll No. \myrollno
\end{flushright}

\clearpage

% ======================================================= ABSTRACT
\thispagestyle{plain}
\addcontentsline{toc}{chapter}{Abstract}
\begin{center}
{\large\bfseries ABSTRACT}
\end{center}
\vspace{0.2in}

Session-based recommendation predicts a user's next interaction from a short
sequence of recent activity, without relying on a persistent user profile.
Large Language Models have recently been applied to this task by expressing the
session and the candidate items in natural language and asking the model to
reason about the user's intent. Frameworks of this kind report substantial gains
over conventional sequential and graph-based recommenders.

These gains are usually reported for a single prompt. A closer reading shows
that the ranking produced by such a system can change when the wording of the
instruction, or the amount of item context supplied, is altered, even though the
session, the candidate set, the model and the decoding configuration remain
identical. Accuracy measured under one prompt therefore does not establish that
the recommender is reliable, and a reproducibility analysis of one such
framework attributes the effect to semantic drift across prompts and domains.

Two mitigations have been proposed in the recent literature. The first optimises
and fuses prompts across domains using a frontier-scale proprietary model. The
second blends the output of a fine-tuned Large Language Model with that of a
separately trained conventional recommender, weighting the two by how sparse a
user's interaction history is. Both are effective, and both assume computational
resources that are unavailable in most academic and practical settings: the
first requires access to a frontier model, and the second requires full-parameter
fine-tuning distributed across a multi-GPU cluster together with a second trained
model.

This work asks whether comparable reliability can be obtained from a single
small, locally hosted model. The proposed method, Reliability-Weighted Rank
Aggregation, scores each candidate pool under several controlled prompt variants
of the same model and combines the resulting rankings using weights derived from
how strongly each variant agrees with the remainder of the ensemble, measured by
Kendall's rank correlation. The same agreement statistic is reported as a
per-session drift score, which expresses semantic drift as an inspectable
quantity rather than inferring it from an aggregate decline in accuracy. A
long-tail-aware extension, adapted from the adaptive aggregation mechanism of a
recent IEEE study, additionally blends the ensemble with the single baseline
prompt according to the length of the session prefix. The method requires no
model training, no second model and no frontier-scale model, and is evaluated on
two domains, MovieLens-1M and Amazon Video Games, under an identical protocol
with paired significance testing.

\vspace{6pt}
\noindent\textbf{Keywords:} session-based recommendation, large language models,
prompt sensitivity, self-consistency, ensemble methods, rank aggregation,
Kendall's tau, semantic drift, reproducibility, cross-domain evaluation.

\clearpage

% ======================================================= CONTENTS
\renewcommand{\contentsname}{TABLE OF CONTENTS}
\tableofcontents
\clearpage

\renewcommand{\listtablename}{LIST OF TABLES}
\addcontentsline{toc}{chapter}{List of Tables}
\listoftables
\clearpage

\pagenumbering{arabic}
\setcounter{page}{1}

% ======================================================= CHAPTER 1
\chapter{Introduction}

\section{Overview}

Recommender systems reduce the burden of choice in environments where the number
of available items far exceeds what a user can inspect. Conventional
collaborative filtering builds a persistent profile for each user from a long
history of interactions. In many practical settings such a profile is either
unavailable or undesirable: the user may not be logged in, the platform may
choose not to retain long-term histories for reasons of privacy, or the user's
interest during the current visit may differ from their historical preference.

Session-based recommendation addresses this setting. It predicts the next item
from a short, ordered sequence of recent interactions alone. The task is
therefore to infer intent from sparse and transient evidence, which makes it
both practically important for anonymous browsing and technically distinct from
profile-based recommendation.

\section{Large Language Models for Session Recommendation}

Early session-based recommenders modelled the interaction sequence with
recurrent networks, and later with graph neural networks and self-attention.
These approaches represent items as learned latent vectors. They capture
collaborative structure effectively, but they operate as opaque mappings from
behaviour to vector space and do not expose an interpretable account of why a
particular item was recommended.

Large Language Models offer a different formulation. The session history and the
candidate items are written out in natural language, and the model is asked to
reason about which candidate the user is most likely to select next. Because the
model brings substantial world knowledge about the items themselves, it can
relate items semantically even when the interaction data are sparse, and it can
be asked to express its reasoning in text. Frameworks following this formulation
report accuracy competitive with, and often superior to, conventional
sequential models.

\section{Limitations of Current Practice}

The shift from latent vectors to natural-language prompting introduces a failure
mode that has no counterpart in conventional recommenders. The output of the
system now depends not only on the interaction data supplied, but on the precise
wording used to express the task and on the quantity of item metadata included
in the prompt.

Three consequences follow. First, an accuracy figure obtained under one prompt
does not establish that the same figure would be obtained under a
semantically equivalent rephrasing, so reported performance is difficult to
interpret and to reproduce. Second, a prompt optimised on one domain may not
transfer to another, since item vocabulary and semantics differ between domains.
Third, the free-form textual output must be parsed into a ranking, and item
names containing numerals can be confused with ranking positions during parsing,
producing outputs that cannot be evaluated at all.

These effects are collectively described in the recent literature as semantic
drift. They are ordinarily observed only indirectly, as a decline in aggregate
accuracy when the domain or the prompt changes.

\section{Motivation}

Two mitigations have been proposed recently, and both are effective within their
own resource assumptions.

The first optimises prompts per domain and then fuses them into a more general
prompt using a frontier-scale proprietary model as a meta-reasoner. This
requires access to that frontier model, which is proprietary, accessed through a
paid interface, and subject to revision or withdrawal by its provider, so the
experiment cannot be pinned to a fixed artefact.

The second augments a conventional recommender with a fine-tuned Large Language
Model, and blends their predictions using a weight derived from the user's
interaction count. This requires full-parameter fine-tuning of a seven-billion
parameter model across sixteen data-centre accelerators, together with the
separate training of a conventional recommender for every dataset considered.

Neither route is available to a researcher whose computational budget is a
single consumer machine. The motivation for this work is the observation that
the instability itself carries usable information. If a recommender produces
different rankings under semantically equivalent prompts, then the degree of
disagreement between those rankings is a measurable signal, obtainable without
labels, without training and without a second model. The present study
investigates whether that signal is sufficient to recover a useful part of the
reliability that the existing mitigations achieve with far greater resources.

\section{Organisation of the Report}

The remainder of this report is organised as follows. Chapter 2 reviews the
relevant literature on session-based recommendation, on Large Language Models
applied to recommendation, and on ensembling and reliability. Chapter 3
identifies the research gaps arising from that review, and Chapter 4 states the
resulting problem. Chapter 5 lists the objectives of the work. Chapter 6
describes the two datasets used and the construction of evaluation instances.
Chapter 7 presents the proposed methodology, including the aggregation
mechanism, the validation strategy and the current implementation status.
Chapter 8 sets out the work plan and schedule, Chapter 9 the expected outcomes,
and Chapter 10 concludes the report.

% ======================================================= CHAPTER 2
\chapter{Literature Survey}

The literature relevant to this work falls into three groups: the development of
session-based recommendation models, the application of Large Language Models to
the recommendation task, and work on reliability, reproducibility and ensembling
that informs the proposed method.

\section{Evolution of Session-Based Recommendation}

Hidasi and colleagues (2016) introduced the use of gated recurrent units for
session-based recommendation, treating the session as a sequence and predicting
the next item from the recurrent hidden state. Subsequent work incorporated
attention to distinguish a user's general interest within a session from their
immediate intent, as in the neural attentive recommendation machine of Li and
colleagues (2017) and the short-term attention and memory priority model of Liu
and colleagues (2018).

Recurrent formulations impose a strict sequential ordering that is not always
appropriate for browsing behaviour. Wu and colleagues (2019) therefore modelled
the session as a graph, allowing non-consecutive transitions between items to be
captured directly. In parallel, self-attention architectures were adopted from
language modelling, first in the self-attentive sequential recommender of Kang
and McAuley (2018) and subsequently in the bidirectional formulation of Sun and
colleagues (2019).

These models are effective at capturing collaborative and sequential structure.
They share a common characteristic relevant to this work: the item is
represented only as a learned vector, so the semantic content of the item, and
the reasoning that connects one item to another, are not represented explicitly.

\section{Large Language Models for Recommendation}

Bao and colleagues (2023) reformulated recommendation as an instruction-following
task and aligned a general-purpose language model to it through instruction
tuning, demonstrating that a comparatively small amount of task-specific tuning
data is sufficient to obtain competitive performance. Broader surveys of the
area, such as that of Zhao and colleagues (2024), document a rapid expansion of
this line of work.

Most directly relevant to the present study, Sun and colleagues (2024) proposed
a prompt-optimisation framework for intent-aware session recommendation. The
framework decomposes the task into intent inference and candidate reranking,
identifies cases in which the ground-truth item was ranked poorly, prompts the
model to reflect on those failures, and iteratively refines the prompt using an
upper-confidence-bound selection strategy. Substantial improvements over
conventional and single-intent baselines are reported across three datasets.

Luo and colleagues (2026) took a complementary approach, combining a fine-tuned
language model with a conventional recommender. Their framework augments the
conventional model with synthetic training data generated by the language model,
augments the language model's prompt with collaborative information drawn from
the conventional model, and blends the two predictions by linear interpolation.
The interpolation weight is a function of the user's interaction count, so that
users with sparse histories are weighted towards the language model. Gains of
roughly fourteen per cent for direct and sequential recommendation are reported.
The authors state explicitly that the fine-tuning procedure, performed on sixteen
data-centre accelerators, is impractical for large-scale deployment.

\section{Reliability, Reproducibility and Ensembling}

Tiwari, Rekha and Mundotiya (2026) conducted a reproducibility analysis of the
prompt-optimisation framework described above. They report that performance
degrades on semantically distinct domains, attributing this to contextual drift
in long sessions and to prompt overfitting arising from a single-domain
optimisation procedure. They further identify a parsing failure mode in which
numerals embedded in item names are interpreted as ranking positions. Their
proposed remedy combines a deterministic index-based output format with a
reflexive cross-domain prompt fusion step driven by a frontier-scale model.

Separately, a substantial body of work in the language-model literature
establishes that sampling a model repeatedly and aggregating the results
improves reliability. Wang and colleagues (2023) introduced self-consistency, in
which multiple reasoning paths are sampled and the most frequently reached answer
is selected. Prompt-ensembling methods extend this idea by combining several
distinct prompts rather than several samples of one prompt. These techniques are
established for question answering and reasoning tasks. They have not, to the
best of the author's knowledge, been evaluated against the specific failure mode
of semantic drift in session recommendation, nor on models small enough to run
on consumer hardware.

The agreement statistic used in this work is due to Kendall (1938), who defined a
rank correlation coefficient based on the proportion of concordant and discordant
pairs between two orderings.

\section{Summary of Reviewed Literature}

Table 2.1 summarises the principal studies reviewed, their focus, their central
idea and the limitations relevant to this work.

\begin{table}[h]
\centering
\small
\caption{Summary of the reviewed literature}
\begin{tabularx}{\textwidth}{@{}p{2.2cm}p{2.6cm}XX@{}}
\toprule
\textbf{Reference} & \textbf{Focus} & \textbf{Key Idea} & \textbf{Limitations}\\
\midrule
{[1]} Hidasi et al. (2016) & Sequential session modelling &
Recurrent network over the interaction sequence &
Items are latent vectors; no semantic reasoning\\[4pt]

{[4]} Wu et al. (2019) & Graph-based sessions &
Session modelled as a graph of item transitions &
Opaque; requires training per dataset\\[4pt]

{[7]} Sun et al. (2024) & LLM prompt optimisation &
Error-driven iterative prompt refinement for intent-aware sessions &
Single static prompt per domain; reliability not measured\\[4pt]

{[8]} Tiwari et al. (2026) & Reproducibility of the above &
Diagnoses semantic drift and parsing failure across domains &
Remedy depends on a frontier-scale proprietary model\\[4pt]

{[9]} Luo et al. (2026) & LLM and conventional model fusion &
Mutual augmentation with interpolation weighted by user sparsity &
Requires fine-tuning on 16 GPUs and a second trained model\\[4pt]

{[10]} Wang et al. (2023) & Self-consistency &
Sample several reasoning paths and take the majority answer &
Demonstrated for reasoning tasks, not recommendation drift\\
\bottomrule
\end{tabularx}
\end{table}

% ======================================================= CHAPTER 3
\chapter{Research Gaps}

The review presented in Chapter 2 exposes three gaps that this work addresses.

\begin{enumerate}[leftmargin=*]
\item \textbf{Reliability under prompt variation is diagnosed but not
quantified per instance.} Existing studies observe semantic drift indirectly, as
a decline in aggregate accuracy when the prompt or the domain changes. No
per-session quantity is reported that would allow an individual recommendation
to be identified as having been produced under unstable conditions. Consequently
a practitioner cannot distinguish a session on which the model was internally
consistent from one on which it was not.

\item \textbf{Existing mitigations assume resources that are unavailable in
most settings.} The two published remedies require, respectively, a
frontier-scale proprietary model to perform cross-domain prompt fusion, and
full-parameter fine-tuning across a multi-GPU cluster together with a second,
separately trained conventional recommender. Neither has been evaluated on a
small, locally hosted, openly licensed model, which is the configuration
available to most academic researchers and to deployments in which data cannot
leave the premises.

\item \textbf{Ensembling has not been connected to this failure mode, and
existing aggregation weights ignore model output.} Self-consistency and
prompt-ensembling methods are established in the language-model literature but
have not been evaluated against semantic drift in session recommendation. Where
an adaptive aggregation weight is used in recommendation, as in the fusion
framework of Luo and colleagues, it is computed from the user's interaction
count alone and is therefore fixed before either model produces an output. No
existing weighting scheme in this area responds to the degree of agreement
between the predictions actually generated.
\end{enumerate}

% ======================================================= CHAPTER 4
\chapter{Problem Statement}

Arising from the gaps identified above, the problem addressed by this work is
stated in three parts.

\begin{enumerate}[leftmargin=*]
\item The sensitivity of a Large Language Model session recommender to prompt
wording and to the quantity of item context has not been measured under
conditions strict enough to attribute the observed variation to those factors
alone, and no per-session measure of that instability is reported.

\item The mitigations proposed for this instability depend on computational
resources that are unavailable in most academic and practical settings, and it
is therefore not known whether comparable reliability can be obtained from a
single small model without any training.

\item Existing aggregation schemes in this area weight their components using
quantities known in advance of inference, and it is not known whether a weight
derived from the agreement between the generated predictions themselves yields a
measurable improvement, nor whether any such improvement generalises across
semantically distinct domains.
\end{enumerate}

Each part of the problem corresponds directly to one of the objectives listed in
Chapter 5.

% ======================================================= CHAPTER 5
\chapter{Objectives}

The objectives of this work are as follows.

\begin{enumerate}[leftmargin=*]
\item To quantify the sensitivity of a Large Language Model session recommender
to prompt wording and to semantic item context under a controlled protocol in
which the session, the candidate pool, the model and the decoding configuration
are held fixed, and to define a per-session measure of ranking instability.

\item To design and implement a reliability-weighted aggregation method that
combines several controlled prompt variants of a single locally hosted model,
without model training, without a second model and without a frontier-scale
model, together with a long-tail-aware extension adapted from the adaptive
aggregation literature.

\item To validate the proposed method against a single-prompt baseline and an
unweighted ensemble on two semantically distinct domains, using paired
statistical testing and an analysis stratified by session history length.
\end{enumerate}

Table 5.1 shows how each objective responds to the corresponding element of the
problem statement.

\begin{table}[h]
\centering
\small
\caption{Mapping of problem statement to objectives}
\begin{tabularx}{\textwidth}{@{}XX@{}}
\toprule
\textbf{Problem} & \textbf{Corresponding Objective}\\
\midrule
Sensitivity is not measured under controlled conditions and no per-session
instability measure exists &
Quantify wording and context sensitivity under a one-factor-at-a-time protocol
and define a per-session drift score\\[6pt]

Published mitigations require a frontier model or a multi-GPU fine-tuning
cluster and a second model &
Design and implement reliability-weighted aggregation over prompt variants of a
single local model, requiring no training\\[6pt]

Aggregation weights ignore the generated predictions, and cross-domain
generality is untested &
Validate against single-prompt and unweighted-ensemble baselines on two domains
with paired significance testing\\
\bottomrule
\end{tabularx}
\end{table}

% ======================================================= CHAPTER 6
\chapter{Dataset Description}

\section{MovieLens-1M}

The MovieLens-1M collection is the primary resource for this work. It contains
1{,}000{,}209 ratings contributed by 6{,}040 users across 3{,}706 distinct
films, together with 3{,}883 rows of film metadata providing the title, the year
of release and a pipe-separated list of genres. The collection is widely used as
a benchmark in the session-recommendation literature and is one of the three
datasets used by the prompt-optimisation framework that motivates this study,
which allows results to be positioned against prior work.

\section{Amazon Video Games (2023)}

The second domain is drawn from the Amazon Reviews 2023 collection released by
the McAuley Laboratory. The Video Games category is used, in the five-core
rating-only form, in which every retained user and every retained item has at
least five interactions. This subset contains 814{,}586 interactions from
94{,}762 users across 25{,}612 items, each interaction recording a user
identifier, an item identifier, a rating and a millisecond-resolution timestamp.
Item metadata providing the product title and its category path is obtained from
the corresponding per-category metadata file.

This domain is chosen deliberately. Product titles in the Video Games category
frequently contain numerals and platform identifiers, which is precisely the
condition under which the parsing failure mode described in the reproducibility
literature arises, and the semantics of the domain differ substantially from
those of film, which makes it a meaningful test of cross-domain generality.

\section{Comparison of the Two Datasets}

Table 6.1 contrasts the characteristics of the two collections.

\begin{table}[h]
\centering
\small
\caption{Comparison of the MovieLens-1M and Amazon Video Games datasets}
\begin{tabularx}{\textwidth}{@{}p{3.6cm}XX@{}}
\toprule
\textbf{Property} & \textbf{MovieLens-1M} & \textbf{Amazon Video Games}\\
\midrule
Interactions & 1{,}000{,}209 & 814{,}586\\
Users (sessions) & 6{,}040 & 94{,}762\\
Distinct items & 3{,}706 & 25{,}612\\
Item context available & Title, year, genres & Title, category path\\
Identifier type & Integer & String (remapped to integer)\\
Timestamp resolution & Seconds & Milliseconds\\
Numerals in item titles & Rare & Frequent\\
Domain noun used in prompt & ``movie'' & ``game''\\
\bottomrule
\end{tabularx}
\end{table}

\section{Construction of Evaluation Instances}

Both collections are reduced to evaluation instances by an identical,
dataset-agnostic procedure, so that no difference in results between the two
domains can be attributed to a difference in preprocessing.

For each user, interactions are sorted by timestamp, with the item identifier
used as a deterministic tie-breaker so that the ordering is reproducible. The
final interaction is withheld as the prediction target and all preceding
interactions form the session prefix presented to the model. Users whose history
is shorter than a specified minimum are excluded. This yields one evaluation
instance per eligible user: 6{,}040 for MovieLens-1M and 94{,}762 for Amazon
Video Games.

Each instance is then assigned a candidate pool of exactly twenty items,
comprising the withheld target and nineteen negative items. Negatives are
selected by popularity computed from training interactions only, with each user's
withheld target removed before the popularity counts are formed, so that the
target cannot influence the construction of its own candidate pool. Items
already present in the session prefix are excluded from the negative set, and
popularity ties are broken by item identifier so that the pool is deterministic.
The position of the target within the pool is cycled systematically across
instances, which prevents the model from benefiting from a fixed target
position.

Because the identifiers used in the Amazon collection are strings, they are
remapped to dense integers through a sorted, deterministic mapping that is
persisted alongside the processed data, so that any synthetic identifier can be
traced back to the original record.

% ======================================================= CHAPTER 7
\chapter{Proposed Methodology}

\section{Overview of the Processing Pipeline}

The proposed pipeline is organised in four stages: preparation of evaluation
instances, generation of ensemble members, aggregation, and evaluation. Table
7.1 summarises the stages and their constituent steps.

\begin{table}[h]
\centering
\small
\caption{Stages of the proposed processing pipeline}
\begin{tabularx}{\textwidth}{@{}p{3.0cm}XX@{}}
\toprule
\textbf{Stage} & \textbf{Steps} & \textbf{Output}\\
\midrule
1. Instance preparation & Load dataset; construct chronological sessions;
build candidate pools & One session and one twenty-item pool per user\\[4pt]

2. Ensemble generation & Render each prompt variant; query the fixed local model;
parse each response strictly & Per-variant candidate scores and ranking\\[4pt]

3. Aggregation & Compute pairwise agreement, reliability weights and the
long-tail weight; combine scores & Aggregated ranking and drift score\\[4pt]

4. Evaluation & Compute accuracy and stability metrics; paired significance
tests; stratify by history length & Per-condition performance and significance\\
\bottomrule
\end{tabularx}
\end{table}

\section{Frozen Experimental Configuration}

All comparisons are made under a single frozen configuration. The model is
Qwen2.5-3B-Instruct, served locally through the Ollama runtime, with the exact
model digest recorded. Decoding is deterministic: temperature is set to zero and
nucleus sampling is disabled. The maximum output length, the response schema and
the parser are identical across every condition.

This configuration is deliberate in one respect that bears on the design of the
method. Because decoding is deterministic, repeating the same prompt produces the
same output, and the diversity required by any ensemble cannot be obtained by
resampling. It must instead be obtained from the prompt itself, which is also the
factor to which the recommender has been shown to be sensitive.

\section{Controlled Prompt Variants}

Two families of controlled variation are defined. Wording variants alter only
the scoring instruction, leaving the session, the candidate list, the output
schema and all surrounding text unchanged. Context variants alter only the item
information presented, ranging from the title alone, through title with genre or
category, to title with genre and year of release.

Each variant carries a stable identifier which is recorded with every trial,
together with the fully rendered prompt, so that any result can be traced to the
exact text that produced it. Unrecognised identifiers are rejected by the
implementation, which prevents an unintended prompt modification from being
recorded as a controlled condition.

\section{Output Protocol and Strict Parsing}

The model is required to return a JSON object containing one numeric score in
the range zero to one hundred for each candidate position. Ranking is then
performed by the program rather than by the model: scores are sorted in
descending order, ties are broken by ascending candidate position, and positions
are mapped back to item identifiers.

This protocol addresses the parsing failure mode reported in the literature. The
model is never asked to reproduce item identifiers, so numerals occurring within
item titles cannot be mistaken for ranking positions. Responses that do not
conform to the schema are recorded as parsing failures together with the raw
text and the reason for rejection; they are never silently repaired, since
repairing them would conceal exactly the unreliability the study sets out to
measure. Parsing failures are reported separately from recommendation failures,
in which a well-formed response simply ranks the target poorly.

\section{Reliability-Weighted Rank Aggregation}

For a given session, let $k$ prompt variants each produce a vector of candidate
scores and a corresponding ranking $r_i$. Agreement between two rankings is
measured by Kendall's rank correlation coefficient, which is $+1$ for identical
orderings, $-1$ for exactly reversed orderings, and zero in expectation for
unrelated orderings.

The mean agreement of member $i$ with the remainder of the ensemble is

\begin{equation}
\tau_i = \frac{1}{k-1}\sum_{j \neq i} \operatorname{KendallTau}(r_i, r_j).
\end{equation}

The reliability weight of member $i$ is obtained by shifting this quantity to be
non-negative and normalising across members,

\begin{equation}
w_i = \frac{\tau_i + 1}{\sum_{j=1}^{k} (\tau_j + 1)},
\end{equation}

and the aggregated score of each candidate is the weighted combination

\begin{equation}
s = \sum_{i=1}^{k} w_i \, s_i .
\end{equation}

The shift by unity is necessary because $\tau_i$ may be negative, and a negative
weight would correspond to inverting a member's ranking rather than discounting
it. Under this formulation a member that disagrees with the remainder of the
ensemble on a particular session receives a reduced weight on that session
alone, and the weight is computed without reference to the ground-truth item, so
the procedure is applicable at inference time.

A member whose response failed to parse is excluded from the ensemble rather
than assigned a default score. If fewer than two members remain, the aggregation
is recorded as a structured failure for that session rather than computed from a
degenerate ensemble.

\section{Per-Session Drift Score}

The same pairwise agreement computation yields a diagnostic quantity. The mean
pairwise Kendall correlation across all members, together with the mean pairwise
Jaccard overlap of the top five and top ten ranked items, is recorded for every
session as a drift score. A session on which the members agree closely receives
a value near unity; a session on which the ensemble divides receives a low or
negative value.

This quantity expresses semantic drift as a directly inspectable per-session
measurement rather than as an inference drawn from an aggregate decline in
accuracy, which addresses the first research gap identified in Chapter 3.

\section{Long-Tail-Aware Weighting}

An extension is adapted from the adaptive aggregation mechanism of Luo and
colleagues (2026). In the original formulation the output of a Large Language
Model is blended with that of a separately trained conventional recommender, and
the blend weight is derived from the logarithm of the user's interaction count,
so that users with sparse histories are weighted towards the language model.

The present work contains only one model family, so the mechanism is repurposed
rather than reused directly. The blend is formed between the aggregated ensemble
and the single baseline prompt, and the weight is derived from the length of the
session prefix. Writing $\ell_s$ for the logarithm of the prefix length
incremented by one, and $\ell_{\min}$ and $\ell_{\max}$ for the corresponding
extremes over the whole dataset,

\begin{equation}
\beta_s = \max\!\left(\frac{\ell_{\max} - \ell_s}{\ell_{\max} - \ell_{\min}},\;
\beta_2\right)\cdot \beta_1,
\end{equation}

\begin{equation}
s_{\text{final}} = \beta_s \, s_{\text{ensemble}} + (1 - \beta_s)\,
s_{\text{baseline}} .
\end{equation}

The parameter $\beta_1$ bounds the weight assigned to the ensemble, and
$\beta_2$ imposes a floor, so that neither component is ever entirely
discarded. The extremes $\ell_{\min}$ and $\ell_{\max}$ are computed once over
the complete dataset rather than over the evaluated subset, so that the weight
assigned to a session does not depend on which sample it is drawn into.

The underlying hypothesis is that a short prefix provides little evidence for
any individual prompt to reason from, so that agreement across prompts is more
informative there, whereas a long prefix may already be handled reliably by a
single prompt. This is stated as a hypothesis to be tested, not as an
established property, and the stratified analysis described below is its test.

\section{Validation Strategy}

Four conditions are compared on identical sessions and identical candidate
pools: the single baseline prompt alone, the unweighted mean of the ensemble,
reliability-weighted aggregation, and the long-tail-weighted blend. Only the
aggregation rule differs between conditions, so any difference in outcome is
attributable to that rule.

Comparisons are restricted to sessions on which both compared conditions
produced a valid ranking. Paired testing is used throughout: a bootstrap
confidence interval is computed over the paired per-session difference in
reciprocal rank, and the Wilcoxon signed-rank test is applied to the same paired
values. All four conditions are additionally recomputed separately on short,
medium and long session-prefix tertiles, which tests directly whether any
improvement is concentrated in sessions with sparse history.

The complete protocol is executed independently on both domains, with no
domain-specific tuning of prompts, weights or parameters, so that cross-domain
generality is assessed rather than assumed.

\section{Evaluation Metrics}

Accuracy is reported as the hit ratio at ranks one, five and ten, and as
normalised discounted cumulative gain at ranks five and ten, computed only over
sessions that produced a valid ranking. Stability is reported as the mean
pairwise Kendall correlation and the mean pairwise Jaccard overlap at ranks five
and ten. Operational reliability is reported as the proportion of parsing
failures and of request failures, which are counted separately from
recommendation errors so that an invalid output is never conflated with a poor
recommendation.

\section{Implementation Status}

The data preparation, prompting, parsing, aggregation and evaluation components
have been implemented as a single reproducible package, accompanied by
seventy-five automated tests which are presently passing. These tests verify the
metric definitions against known values, confirm that the aggregation
down-weights a synthetically divergent member, and confirm that parsing failures
and degenerate ensembles are reported rather than silently repaired. The complete
pipeline has been executed end to end on both domains using a deterministic
substitute model, which validates the data path without consuming inference time.

Evaluation using the real model has so far been limited to a small pilot on
MovieLens-1M, conducted to confirm that the local model produces parseable output
under each wording variant. The pilot comprised five instances per condition and
is reported in Chapter 9 solely as evidence that the instability under
investigation is present; it is not treated as a measurement of performance. The
full-scale evaluation described in Section 7.7 remains to be executed and is
scheduled in Chapter 8.

% ======================================================= CHAPTER 8
\chapter{Work Plan}

The work is scheduled over twelve months and divided into three phases. Phase 1
covers the study, the experimental protocol and data preparation. Phase 2 covers
the design and implementation of the aggregation method. Phase 3 covers
full-scale evaluation, statistical validation and documentation. Table 8.1 lists
the work packages and their planned durations.

\begin{table}[h]
\centering
\small
\caption{Work plan and schedule}
\begin{tabularx}{\textwidth}{@{}p{0.6cm}Xp{1.8cm}p{2.0cm}@{}}
\toprule
\textbf{No.} & \textbf{Task / Work Package} & \textbf{Phase} & \textbf{Duration}\\
\midrule
1 & Literature survey and research gap identification & Phase 1 & M1 -- M2\\
2 & Experimental protocol and frozen configuration & Phase 1 & M2 -- M3\\
3 & Dataset acquisition and session construction & Phase 1 & M3 -- M4\\
4 & Candidate pool generation and leakage controls & Phase 1 & M4 -- M5\\
5 & Prompt variant design and strict output parser & Phase 2 & M5 -- M6\\
6 & Reliability-weighted aggregation and drift score & Phase 2 & M6 -- M7\\
7 & Long-tail-aware weighting extension & Phase 2 & M7 -- M8\\
8 & Cross-domain extension to the second dataset & Phase 2 & M8 -- M9\\
9 & Full-scale evaluation on both domains & Phase 3 & M9 -- M10\\
10 & Statistical and stratified analysis & Phase 3 & M10 -- M11\\
11 & Comparison with published baselines & Phase 3 & M11\\
12 & Thesis writing and paper submission & Phase 3 & M11 -- M12\\
\bottomrule
\end{tabularx}
\end{table}

The phases are deliberately overlapped at their boundaries, so that the
aggregation method is implemented while the prompt variants are still being
refined, and statistical analysis begins while the full-scale evaluation is
still producing results. This reduces the risk that a delay in one package
propagates through the whole schedule. Packages 1 to 8 are complete at the time
of writing; packages 9 to 12 are pending the availability of accelerated
computing resources.

% ======================================================= CHAPTER 9
\chapter{Expected Outcomes}

The following outcomes are anticipated from this work.

\begin{itemize}[leftmargin=*]
\item A quantitative characterisation of how far the ranking produced by a Large
Language Model session recommender moves when only the prompt wording or the
item context is altered, obtained under a protocol in which every other factor
is held fixed. A preliminary pilot on five instances already exhibits variation
in the hit ratio at rank one across wording variants, which motivates the
full-scale measurement but does not constitute a result.

\item A per-session drift score that expresses semantic drift as an inspectable
quantity, allowing an individual recommendation to be identified as having been
produced under unstable conditions rather than requiring instability to be
inferred from aggregate accuracy.

\item An aggregation method that improves reliability using a single locally
hosted model, requiring no model training, no second model and no frontier-scale
model, evaluated against both a single-prompt baseline and an unweighted
ensemble under paired significance testing.

\item Evidence on whether the benefit of ensembling is concentrated in sessions
with short interaction histories, obtained from the stratified analysis. A
negative result here is itself informative, since it would indicate that prompt
sensitivity is not governed by the quantity of available history.

\item Evidence on whether the method transfers between semantically distinct
domains without domain-specific tuning, obtained by executing an identical
protocol on film and video-game data.

\item A reproducible open implementation, with the model digest, decoding
configuration, prompt text, candidate pools and raw responses recorded for every
trial, so that the reported figures can be regenerated independently.
\end{itemize}

% ======================================================= CHAPTER 10
\chapter{Conclusion}

Large Language Models have been applied to session-based recommendation with
encouraging results, but the reported accuracy is generally obtained under a
single prompt, and the ranking produced by such a system has been shown to move
when that prompt is rephrased or when the item context is altered. Existing
remedies for this instability are effective but assume either a frontier-scale
proprietary model or full-parameter fine-tuning across a multi-GPU cluster
alongside a second trained model, and neither assumption holds in most academic
or privacy-constrained settings.

This work proceeds from the observation that the instability itself is a usable
signal. If semantically equivalent prompts produce different rankings, the
degree of disagreement between those rankings can be measured without labels,
without training and without a second model. The proposed method scores each
candidate pool under several controlled prompt variants of one small, locally
hosted model, weights each variant by its agreement with the remainder of the
ensemble, and reports the same agreement statistic as a per-session drift score.
A long-tail-aware extension, adapted from a recent adaptive aggregation
mechanism and repurposed for a single-model setting, blends the ensemble with the
single baseline prompt according to the length of the session prefix.

The design is deliberately conservative in what it claims. Agreement between
prompt variants is a proxy for reliability and not for correctness: if every
member of the ensemble shares the same error, the method will weight them
equally and return that error confidently. The approach also forgoes the
collaborative signal that a conventional recommender contributes in the fusion
framework from which the weighting mechanism is adapted, and it incurs a
multiple of the inference cost of a single prompt. What it offers in exchange is
reliability obtained on hardware that is widely available, under a configuration
that can be pinned exactly and reproduced by others.

The value of the outcome accordingly lies less in obtaining the highest figure
than in obtaining a figure that can be trusted, reproduced and inspected at the
level of the individual session. Should the method prove effective, the practical
implication is that improved reliability in prompt-based recommendation need not
be purchased with a frontier model or a training cluster, but can be obtained by
asking a small model the same question in several ways and listening to how
consistently it answers.

% ======================================================= REFERENCES
\chapter*{REFERENCES}
\addcontentsline{toc}{chapter}{References}
\markboth{REFERENCES}{REFERENCES}

\begin{enumerate}[label={[\arabic*]},leftmargin=*,itemsep=6pt]

\item Hidasi, B., Karatzoglou, A., Baltrunas, L., \& Tikk, D. (2016).
Session-based recommendations with recurrent neural networks. \emph{International
Conference on Learning Representations (ICLR)}.

\item Li, J., Ren, P., Chen, Z., Ren, Z., Lian, T., \& Ma, J. (2017). Neural
attentive session-based recommendation. \emph{Proceedings of the 26th ACM
International Conference on Information and Knowledge Management (CIKM)},
1419--1428.

\item Liu, Q., Zeng, Y., Mokhosi, R., \& Zhang, H. (2018). STAMP: Short-term
attention/memory priority model for session-based recommendation.
\emph{Proceedings of the 24th ACM SIGKDD International Conference on Knowledge
Discovery and Data Mining}, 1831--1839.

\item Wu, S., Tang, Y., Zhu, Y., Wang, L., Xie, X., \& Tan, T. (2019).
Session-based recommendation with graph neural networks. \emph{Proceedings of the
AAAI Conference on Artificial Intelligence}, 33(1), 346--353.

\item Kang, W.-C., \& McAuley, J. (2018). Self-attentive sequential
recommendation. \emph{IEEE International Conference on Data Mining (ICDM)},
197--206.

\item Sun, F., Liu, J., Wu, J., Pei, C., Lin, X., Ou, W., \& Jiang, P. (2019).
BERT4Rec: Sequential recommendation with bidirectional encoder representations
from transformer. \emph{Proceedings of the 28th ACM International Conference on
Information and Knowledge Management (CIKM)}, 1441--1450.

\item Sun, Z., Liu, H., Qu, X., Feng, K., Wang, Y., \& Ong, Y. S. (2024). Large
language models for intent-driven session recommendations. \emph{Proceedings of
the 47th International ACM SIGIR Conference on Research and Development in
Information Retrieval}, 324--334.

\item Tiwari, A., Rekha, K. N. L., \& Mundotiya, R. K. (2026). A reproducibility
analysis of PO4ISR: Diagnosing and mitigating semantic drift in LLM-based session
recommendation. \emph{arXiv preprint arXiv:2605.18780}.

\item Luo, S., Yao, Y., He, B., Shao, W., Xu, J., Huang, Y., Zhou, A., Zhang, X.,
Xiao, Y., Hou, H., Zhan, M., \& Song, L. (2026). Integrating large language
models into recommendation via mutual augmentation and adaptive aggregation.
\emph{IEEE Journal of Selected Topics in Signal Processing}, 20(1), 77--89.

\item Wang, X., Wei, J., Schuurmans, D., Le, Q. V., Chi, E. H., Narang, S.,
Chowdhery, A., \& Zhou, D. (2023). Self-consistency improves chain of thought
reasoning in language models. \emph{International Conference on Learning
Representations (ICLR)}.

\item Bao, K., Zhang, J., Zhang, Y., Wang, W., Feng, F., \& He, X. (2023).
TALLRec: An effective and efficient tuning framework to align large language
model with recommendation. \emph{Proceedings of the 17th ACM Conference on
Recommender Systems}, 1007--1014.

\item Zhao, Z., Fan, W., Li, J., Liu, Y., Mei, X., Wang, Y., Wen, Z., Wang, F.,
Zhao, X., Tang, J., \& Li, Q. (2024). Recommender systems in the era of large
language models (LLMs). \emph{IEEE Transactions on Knowledge and Data
Engineering}, 36(11), 6889--6907.

\item Kendall, M. G. (1938). A new measure of rank correlation.
\emph{Biometrika}, 30(1--2), 81--93.

\item Harper, F. M., \& Konstan, J. A. (2015). The MovieLens datasets: History
and context. \emph{ACM Transactions on Interactive Intelligent Systems}, 5(4),
1--19.

\item Hou, Y., Li, J., He, Z., Yan, A., Chen, X., \& McAuley, J. (2024). Bridging
language and items for retrieval and recommendation. \emph{arXiv preprint
arXiv:2403.03952}.

\end{enumerate}

\end{document}
